\exercice{Molécules organiques (5 minutes) }\points{4}

\question{
  Justifier que \variationSujet{le carbone}{l'oxygène}, de numéro atomique $Z =$ \variationSujet{6}{8}, forme \variationSujet{4}{2} liaisons covalentes, en utilisant sa configuration électronique.\points{2}
}{}

\begin{center}
  \variationSujet{\chemfig{!\alanineSemiDev}}{\chemfig{H-N(-[:-90]H) - C (-[:-90]H) (-[:90] H) -C (!\lldh O) !\ldb O - H}}

  \legende{Molécule \variationSujet{d'alanine}{de glycine}}
\end{center}

\question{
  Donner le nom de la représentation (développée, semi-développée, topologique) de la molécule \variationSujet{d'alanine}{de glycine}.\points{1}
}{}

\question{
  Donner la formule brute de la molécule \variationSujet{d'alanine}{de glycine}.\points{1}
}{}

% \question{
%   Donner la formule développée de la molécule d'alanine.
% }{}[6]